\chapter{X 射线}
\problem{8.1}{某 X 光机的高压为 10 万伏，问发射光子的最大能量多大？算出发射 X 光的最短波长。}
\answerarea
\problem{8.2}{利用普通光学反射光栅可以测定 X 光波长。当掠射角为 $\theta$ 而出现 $n$ 级极大之出射光线偏离入射光线为 $2\theta+\alpha$（见习题图 8.1，$\alpha$ 为偏离 0 级极大出射线的角度）。试证：出现 $n$ 级极大的条件是
\[
2d\sin\frac{2\theta+\alpha}{2}\sin\frac{\alpha}{2}=n\lambda,
\]
$d$ 为光栅常数（即两刻纹中心之间的距离）。当 $\theta$ 和 $\alpha$ 都很小时公式简化为
\[
d\left(\theta\alpha+\frac{\alpha^2}{2}\right)=n\lambda.
\]
\fig{8_1}{8.1}}
\answerarea\clearpage
\problem{8.3}{一束 X 光射向每毫米刻有 100 条纹的反射光栅，其掠射角为 $20'$。已知第一级极大出现在离 0 级极大出射线的夹角也是 $20'$。算出入射 X 光的波长。}
\answerarea
\problem{8.4}{已知 $\mathrm{Cu}$ 的 $\mathrm K_\alpha$ 线波长是 $1.542\,\angstrom$，以此 X 射线与 $\mathrm{NaCl}$ 晶体自然面成 $15^\circ50'$ 角入射而得第二级极大，试求 $\mathrm{NaCl}$ 晶体常数 $d$。}
\answerarea\clearpage
\problem{8.5}{铝（$\mathrm{Al}$）被高速电子束轰击而产生的连续 X 光谱的短波限为 $5\,\angstrom$。问这时是否也能观察到其标识谱 $\mathrm K$ 系线？}
\answerarea
\problem{8.6}{已知 $\mathrm{Al}$ 和 $\mathrm{Cu}$ 对于 $\lambda=0.7\,\angstrom$ 的 X 光的质量吸收系数分别是 $0.5\,\mathrm{m^2/kg}$ 和 $5.0\,\mathrm{m^2/kg}$，$\mathrm{Al}$ 和 $\mathrm{Cu}$ 的密度分别是 $2.7\times10^3\,\mathrm{kg/m^3}$ 和 $8.93\times10^3\,\mathrm{kg/m^3}$。现若分别单独用 $\mathrm{Al}$ 板或 $\mathrm{Cu}$ 板作挡板，要使波长为 $0.7\,\angstrom$ 的 X 光的强度减至原来强度的 $1/100$，问要选用的 $\mathrm{Al}$ 板或 $\mathrm{Cu}$ 板应多厚？}
\answerarea\clearpage
\problem{8.7}{为什么在 X 光吸收光谱中 $\mathrm K$ 系带的边缘是简单的，$\mathrm L$ 系带是三重的，$\mathrm M$ 系带是五重的？}
\answerarea
\problem{8.8}{试证明 X 光标识谱和碱金属原子光谱有相仿的结构。}
\answerarea\clearpage
